\documentclass[11pt]{article}
\usepackage[spanish]{babel}
\usepackage[T1]{fontenc}
\usepackage{amsmath,amsfonts,amssymb,amsthm}
\usepackage{amsfonts}
\usepackage{graphicx}
\usepackage{amssymb}
\usepackage{geometry}
\usepackage{float}
\usepackage{fourier,libertine}
\usepackage{xcolor}
\usepackage{tikz}
\usepackage{geometry}
\usepackage{picinpar}
\usepackage{accents}
\newcommand{\interior}[1]{\accentset{\smash{\raisebox{-0.12ex}{$\scriptstyle\circ$}}}{#1}\rule{0pt}{2.3ex}}
\fboxrule0.0001pt \fboxsep0pt
 \geometry{
 a4paper,
 total={170mm,245mm},
 left=20mm,
 top=30mm,
 }
\pagestyle{empty}
\begin{document}
\setlength{\parindent}{0cm}
\hoffset-0.46cm
\voffset-1.46cm
\begin{window}[0,l,{\includegraphics[scale=0.22]{Graphics/logo.eps}},]
\large\scshape \hspace{1.4cm}\textsf{Universidad Nacional de Colombia} \\
\textcolor{white}{\tiny.}  \large \hspace{2.4cm} \textsf{Facultad de Ciencias} \\
\textcolor{white}{\tiny.}   \normalsize\hspace{2.8cm}\textsf{Topología General}\\
\hspace*{3.9cm}\textsf{Taller II}\\
\end{window}
\vspace{0.8cm}
\textsf{Bourbaki} 
\normalsize
\dotfill
\vspace{0.7cm}


\begin{enumerate}
    \item Siendo $a$, $b$, $c$, $z_0$ constantes complejas, demostrar:
    \[
    \lim_{z \to z_0} c = c; \quad \lim_{z \to z_0} (a z + b) = a z_0 + b; \quad \lim_{z \to z_0} (z^2 + c) = z_0^2 + c
    \]
    \[
    \lim_{z \to z_0} \operatorname{Re}(z) = \operatorname{Re}(z_0); \quad \lim_{z \to z_0} \overline{z} = \overline{z_0}; \quad \lim_{z \to 0} \frac{\overline{z}^2}{z} = 0
    \]

\textbf{Solucao:} $\lim_{z \to z_0} c = c$, note que dado $\epsilon>0$, tomando $\delta=$lo que le de la gana

$$|f(z)-f(z_0)|=|c-c|=0<\epsilon$$

$\lim_{z\to z_0}(az+b)=az_0+b$, eso es trivial allí. Dado $\epsilon>0$, sea $\delta=\frac{\epsilon}{|a|}$, note que si $|z-z_0|<\delta$ entonces

\begin{align*}
  |f(z)-f(z_0)|=|az+b-az_0-b|=|a||z-z_0|<|a|\delta=\epsilon
.\end{align*}

Luego viene uno que usted tiene que hacer cositas pero pues no deja de ser una maricada, $\lim_{z\to z_0}(z^2+c)=z_0^2+c$. Usted primero agarre $\delta_1\leq 1$, note que si

$$|z|-|z_0|\leq|z-z_0|<\delta\leq1$$

entonces 

$$|z+z_0|\leq|z|+|z_0|<1+2|z_0|$$

entonces ponga $\delta_2=\frac{\epsilon}{1+2|z_0|}$ y su delta es $\delta=\min\{\delta_1,\delta_2\}$, y eso le sirve para todo epsilon porque pille:

$$|f(z)-f(z_0)|=|z^2-z_0^2|=|z-z_0||z+z_0|<\epsilon$$

y por eso es que uno debe parar bolas en diferencial.\\

$\lim_{z\to z_0}\Re(z)=\Re(z_0)$, sea $\delta=\sqrt{\epsilon}$, note que para todo $\epsilon>0$

$$|\Re(z)-\Re(z_0)|=|\Re(z-z_0)|\leq|z-z_0|^2<\delta^2=\epsilon$$\\

$\lim_{z \to z_0} \overline{z}=\overline{z_0}$, sea $\delta=\epsilon$, para todo $\epsilon>0$, tenemos que:

$$|\overline{z}-\overline{z_0}|=|\overline{z-z_0}|=|z-z_0|<\epsilon.$$

Para el último tome $\delta=\epsilon$, note que para todo $\epsilon>0$ tenemos que

$$\left|\frac{\overline{z}^2}{z} \right|=\frac{|\overline{z}||\overline{z}|}{|z|}=|\overline{z}|=|z|\leq\epsilon.$$

    \item Calcular:
    \[
    \lim_{z \to \infty} \frac{4z^2}{(z - 1)^2}; \quad \lim_{z \to 1} \frac{1}{(z - 1)^3}; \quad \lim_{z \to \infty} \frac{z^2 + 1}{z - 1}
    \]

    Ahí dice calcular, entonces eso da 4, infinito y el otro también, si no me cree pues hagalo ud

    \item Demostrar que si $f$ es continua y no nula en un punto $z_0$, entonces $f(z) \neq 0$ para todo $z$ en alguna vecindad de $z_0$ $(V_{z_0})$.

\begin{proof}
Suponga que no, ahí van a pasar cosas, lo que pasa allí es que existe $z\in B(z_0,\delta)$ tal que $f(z)=0$ tome $\epsilon=|f(z_0)|$, como en esa bola la función es continua entonces

$$|f(z)-f(z_0)|=|f(z_0)|<\epsilon=|f(z_0)|$$
si pilla que le contradice cosas?
\end{proof}

    \item Sea $A \subset \mathbb{C}$ un conjunto compacto. Mostrar que si $f$ es una función continua definida sobre $A$, entonces $f(A)$ es compacto.

\begin{proof}
Sea $\{z_n\}$ una sucesión en $A$, como $A$ es compacto entonces tiene una subsucesión convergente, digamos $\{z_{k}\}$ que converge a $z$, esto que dado $\delta>0$, existe un $K$ tal que si $k>K$ entonces $|z_{k}-z|<\delta$, por la continuidad $|f(z_k)-f(z)|<\epsilon$ para todo $\epsilon>0$. Acabó porque eso le dice que allá la sucesión $f(z_k)$ converge.
\end{proof}

    \item Sea $A \subset \mathbb{C}$ un conjunto compacto, y $f: A \to \mathbb{R}$ una función continua. Demostrar que $f$ tiene un máximo sobre $A$.
    \begin{proof}
    Facilito, como $f$ es contínua y $A$ es compacto, $f[A]$ es compacto en $\mathbb{R}$, y como en $\mathbb{R}$ todo conjunto compacto es cerrado y acotado, pues por la completez, $f[A]$ tiene supremo, además $\sup\{f[A]\}\in f[A]$ pues $f[A]$ es cerrado, en otras palabras, $f[A]$ tiene máximo.
    \end{proof}
    \item Sea $A \subset \mathbb{C}$ un conjunto compacto, y $f: A \to \mathbb{C}$ una función continua. Demostrar que $f$ es uniformemente continua.
    \begin{proof}
    Suponga que no, luego existe $\epsilon>0$ tal que si $\delta_n=\dfrac{1}{n}, n\in\mathbb{N}$ existen $z_n,w_n\in A$ tales que $|z_n-w_n|<\dfrac{1}{n}$  pero $|f(z_n)-f(w_n)|>\epsilon$.

    Como $A$ es compacto y $\{z_n\}\subseteq A$, existe una subsucesión $\{z_i\}_{i\in I}$ con $I\subseteq\mathbb{N}$ tal que dado $\epsilon$, existe $N \in I$ tal que si $i > N$ entonces $|z_i-z|<\dfrac{\epsilon}{3}$ para algún $z\in A$. 

    Ya que a cada $z_n$ le corresponde un $w_n$, tenemos la subsucesión $\{w_i\}_{i\in I} \subseteq A$ y de nuevo por la compacidad de $A$, obtenemos una subsucesión $\{w_j\}_{j\in J}$ con $J\subseteq I$ de manera que existe $M \in J$ tal que si $j > M$ entonces $|w_j-w|<\dfrac{\epsilon}{3}$. También tenemos que existe $L\in J$ tal que si $j>L$, entonces $\dfrac{1}{j}<\dfrac{\epsilon}{3}$.

    Aplicando desigualdad triangular dos veces y tomando $K=\max\{N,M,L\}$, tenemos que si $j>K$, entonces:
    \[
    |z-w|\leq |z-z_j|+|z_j - w_j|+|w_j - w| \leq \dfrac{\epsilon}{3}+\dfrac{\epsilon}{3}+\dfrac{\epsilon}{3}=\epsilon.
    \]
    Así, tenemos que $z=w$.

    Como $f$ es contínua, tenemos que si $j>K$, $|f(z_j)-f(z)|<\dfrac{\epsilon}{2}$ y $|f(w_j)-f(w)|<\dfrac{\epsilon}{2}$, en conclusión:
    \[
    |f(z_j)-f(w_j)|\leq|f(z_j) - f(z)|+|f(z)-f(w)|+ |f(w)-f(w_j)|=|f(z_j) - f(z)|+|f(w)-f(w_j)|< \dfrac{\epsilon}{2}+ \dfrac{\epsilon}{2}=\epsilon
    \]
    Lo cual contradice que $|f(z_n)-f(z_w)|>\epsilon$.

    \end{proof}
    Esta prueba es una porquería pero así tocaba hacerla entonces bueno, fecho.
\end{enumerate}
\end{document}